\scanpage{118}
\[ \alpha,\beta,\gamma\in[-1,2]\qquad\alpha\leq\beta\leq\gamma \]

\[
\begin{aligned}
c_{\beta,\gamma}&=\text{Nb de structures }\unsure{\mathcal T|\mathscr T|\tau}_\beta\ \text{dans un }\mathcal T_\gamma\\
c'_{\beta\gamma}&=\rule{3.1cm}{.4pt}\ \text{épinglées }\mathcal T_\beta\ \text{dans un }\mathcal T_\gamma\\
N_\beta&=\inserted{c'_{\beta\beta}=}\text{ordre grpe autom de }\mathcal T_\beta
\end{aligned}
\]

\[ \boxed{c'_{\beta\gamma}=N_\beta c_{\beta\gamma}}\tag*{(1)} \]

\deleted{\unsure{NB Noter aussi|$N_\beta$} $N_{\unsure{\alpha|2}}=c'_{\alpha\alpha}$} i.e.

\[ \boxed{c_{\beta\gamma}=\frac{c'_{\beta\gamma}}{c'_{\beta\beta}}} \]

\[ \boxed{c'_{\beta\gamma}=c'_{\alpha,\gamma}\unsure{\nu|V}_{\alpha,\beta,\gamma}}\tag*{(2)} \]

où $\nu_{\alpha\beta\gamma}={}$ nb des structures épinglées $\mathcal T_\beta$ dans un $\mathcal T_\gamma$, prolongeant une structure épinglée $\mathcal T_\alpha$

On calcule les \deleted{\illegible{1}} $c_{\beta\gamma},c'_{\beta\gamma}$ ($\unsure{\gamma|\beta}$ fixé, \inserted{tous les} $\unsure{\alpha|\beta}\leq\unsure{\beta|\gamma}$) par récurrence sur $\beta$, \deleted{Pour $\gamma$ fixé on récurre} \deleted{\unsure{Comme ça fait} le total de $c'_{\beta\gamma}$ et \unsure{groupe}} \inserted{Si $\beta<\gamma$} \unsure{supposant|supposons} que les $c'_{\alpha\beta}$ sont connus pour $\alpha\leq\beta<\gamma$). Alors \unsure{cl|calcul} de $c_{\beta\gamma}$ \deleted{\unsure{et}} \unsure{soit} \deleted{\unsure{il suffit de}} \unsure{déduira|déduire} de \unsure{celui} de $c'_{\beta\gamma}$ \deleted{\unsure{via}} \unsure{connu} déjà $N_\beta$ et

\[ c'_{\beta\gamma}=N_\beta c_{\beta\gamma} \]

Si $\beta=\gamma$, $c_{\beta\gamma}=1$, \deleted{$c'_{\beta\gamma}={}$} il reste à connaître $c'_{\beta\gamma}$. \deleted{et choisissons} et pour $\beta=-1$ ($\mathcal T_\beta$ triangle) on a formules

\[
\begin{aligned}
c_{-1,\gamma}&=\text{nb de triangles }\inserted{\text{de }\mathcal T_\gamma}=(1+c_\gamma)(1+2c_\gamma)\\
c'_{-1,\gamma}&=\text{nb de triangles épinglés de }\mathcal T_\gamma=6(1+c_\gamma)(1+2c_\gamma)
\end{aligned}
\]

On \unsure{récurre} sur $\beta$ \unsure{$=-1$}. Si $\beta\unsure{\neq|>}-1$ (donc $\gamma\geq0$, i.e. $\gamma=0,1$ ou $2$) soit $\alpha<\beta$ \inserted{(on prend $\alpha=\beta-1$)} \deleted{On a de} alors la formule (2) nous donne $c'_{\beta\gamma}$ en fonction de $c'_{\alpha\gamma}$ (déjà connu) et $\nu_{\alpha\beta\gamma}$ (qu'il faut donc déterminer).

\deleted{Reste (en plus des $\nu_{\alpha\beta\gamma}$) à déterminer $N_{\unsure{\gamma|\beta}}$. Or \deleted{comme} pour $\beta=\gamma-1$, \unsure{comme|considérer} le grpe $G_\gamma$ opérant sur l'ens des $\mathcal T_\beta$ épinglés plongés dans $\mathcal T_\gamma$, (au nombre $c'_{\beta\gamma}$) alors le stabilisateur $G_{\beta\gamma}\subset\unsure{G_\gamma}$ d'un él. est d'indice $c'_{\beta\gamma}$, et son ordre est l'ordre $N_{\beta\gamma}$ du grpe des}

% archive sheet number: 117
